%% toc: P15 169 169.9,169.10,169.11
% Proofs for the large-sample regimes of Student inference.
\section{Large-sample regimes of Student inference: proofs}
\label{sec:annex-p15-w02b-p15-059}

Notation is that of \cref{def:p15-w02b-p15-059-setting}. Throughout,
\(G_t(x)=\Phi(t\sqrt x)\) for real \(t\) and \(x>0\). Since
\(\phi_0'(y)=-y\phi_0(y)\), direct differentiation gives, for \(x>0\),
\[
 G_t'(x)=\frac{t\,\phi_0(t\sqrt x)}{2\sqrt x},\qquad
 G_t''(x)=-\frac{t\,\phi_0(t\sqrt x)\,(1+t^2x)}{4x^{3/2}} .
\]
Indeed, differentiating \(G_t'\) gives
\[
 G_t''(x)=\frac t2\Bigl[-t\sqrt x\,\phi_0(t\sqrt x)\cdot\frac t{2\sqrt x}\cdot
 x^{-1/2}-\frac12\phi_0(t\sqrt x)x^{-3/2}\Bigr],
\]
which simplifies to the displayed \(G_t''\). Both are jointly continuous in \((t,x)\in\mathbb R\times(0,\infty)\),
and \(G_t''(1)=-\psi(t)\) with \(\psi\) as in
\cref{thm:p15-w02b-p15-059-second-order}.

\begin{annexitem}{annex:p15-w02b-p15-059-chisq-moments}{thm:p15-w02b-p15-059-chisq-moments}{Chi-square moments}
\begin{proof}
Put \(a=\nu/2\). By \cref{def:p15-gamma-density} the law \(\chi^2_\nu\) has
density \(2^{-a}x^{a-1}e^{-x/2}/\Gamma(a)\) on \((0,\infty)\); being
absolutely continuous it has no atoms, and it gives mass one to
\((0,\infty)\), so \(W>0\) almost surely. Substituting \(x=2y\),
\[
 EW^k=\frac{2^{-a}}{\Gamma(a)}\int_0^\infty x^{a+k-1}e^{-x/2}\,dx
     =\frac{2^k\,\Gamma(a+k)}{\Gamma(a)} .
\]
For \(b>0\), integration by parts on \([\eta,R]\) and the limits
\(\eta\downarrow0\), \(R\to\infty\) give
\(\Gamma(b+1)=\int_0^\infty y^be^{-y}\,dy=b\int_0^\infty y^{b-1}e^{-y}\,dy
=b\,\Gamma(b)\), because the boundary terms \(\eta^be^{-\eta}\) and
\(R^be^{-R}\) tend to zero. Hence
\(\Gamma(a+k)/\Gamma(a)=a(a+1)\cdots(a+k-1)\) and
\(EW^k=\prod_{j=0}^{k-1}(\nu+2j)\).

Write \(m_k=EW^k\). Then \(m_1=\nu\) and
\(E(W-\nu)^2=m_2-\nu^2=\nu(\nu+2)-\nu^2=2\nu\). For the fourth central
moment, binomial expansion gives
\(E(W-\nu)^4=m_4-4\nu m_3+6\nu^2m_2-4\nu^3m_1+\nu^4
=m_4-4\nu m_3+6\nu^2m_2-3\nu^4\), with
\[
\begin{aligned}
 m_4&=\nu^4+12\nu^3+44\nu^2+48\nu,\\
 4\nu m_3&=4\nu^2(\nu^2+6\nu+8)=4\nu^4+24\nu^3+32\nu^2,\\
 6\nu^2m_2&=6\nu^4+12\nu^3 .
\end{aligned}
\]
Collecting powers, the coefficients of \(\nu^4\) and \(\nu^3\) are
\(1-4+6-3=0\) and \(12-24+12=0\), and what remains is
\(12\nu^2+48\nu\).
\end{proof}
\end{annexitem}

\begin{annexitem}{annex:p15-w02b-p15-059-student-limit}{thm:p15-w02b-p15-059-student-limit}{Density, CDF and quantile limits}
\begin{proof}
\emph{Part 1.} Put \(g_\nu(t)=(1+t^2/\nu)^{-(\nu+1)/2}\). By
\cref{thm:p15-student-law}, \(f_\nu=K_\nu g_\nu\) for a constant
\(K_\nu>0\), and \(\int f_\nu\,d\lambda=1\) gives
\(K_\nu=1/\int g_\nu\,d\lambda\).

\emph{Pointwise limit of \(g_\nu\).} For \(u\geq0\),
\(u-u^2/2\leq\log(1+u)\leq u\): both differences vanish at \(u=0\) and
have derivatives \(1-1/(1+u)\geq0\) and \(1/(1+u)-1+u=u^2/(1+u)\geq0\).
With \(u=t^2/\nu\),
\[
 \frac{\nu+1}2\Bigl(\frac{t^2}\nu-\frac{t^4}{2\nu^2}\Bigr)
 \leq-\log g_\nu(t)\leq\frac{\nu+1}2\cdot\frac{t^2}\nu ,
\]
and both bounds tend to \(t^2/2\). Hence
\(g_\nu(t)\to e^{-t^2/2}\).

\emph{Domination.} For \(x\geq0\) the function
\(r\mapsto r\log(1+x/r)\) is nondecreasing on \(r>0\): its derivative is
\(\log(1+u)-u/(1+u)\) with \(u=x/r\), which vanishes at \(u=0\) and has
derivative \(u/(1+u)^2\geq0\) in \(u\). Hence for \(\nu\geq2\),
\(\nu\log(1+t^2/\nu)\geq2\log(1+t^2/2)\), that is
\((1+t^2/\nu)^{\nu/2}\geq1+t^2/2\). Since also
\((1+t^2/\nu)^{-1/2}\leq1\),
\[
 0<g_\nu(t)\leq(1+t^2/\nu)^{-\nu/2}\leq\frac1{1+t^2/2}\qquad(\nu\geq2),
\]
and \(\int(1+t^2/2)^{-1}\,dt=\sqrt2\,\pi<\infty\). Dominated convergence
and \cref{thm:p15-normal-normalization} give
\(\int g_\nu\,d\lambda\to\int e^{-t^2/2}\,d\lambda=\sqrt{2\pi}\). Hence
\(K_\nu\to1/\sqrt{2\pi}\) and \(f_\nu(t)=K_\nu g_\nu(t)\to\phi_0(t)\).

\emph{Part 2.} Both \(f_\nu\) and \(\phi_0\) are probability densities and
\(f_\nu\to\phi_0\) everywhere, so \cref{thm:p15-scheffe-density-limit} gives
\(\int|f_\nu-\phi_0|\,d\lambda\to0\). Explicitly:
\(0\leq(\phi_0-f_\nu)^+\leq\phi_0\) and \((\phi_0-f_\nu)^+\to0\), so
\(\int(\phi_0-f_\nu)^+\to0\) by dominated convergence; and
\(\int(\phi_0-f_\nu)=0\) gives
\(\int|f_\nu-\phi_0|=2\int(\phi_0-f_\nu)^+\). For every real \(t\),
\[
 |F_\nu(t)-\Phi(t)|=\Bigl|\int_{-\infty}^t(f_\nu-\phi_0)\,d\lambda\Bigr|
 \leq\int_{\mathbb R}|f_\nu-\phi_0|\,d\lambda ,
\]
and the bound does not depend on \(t\).

\emph{Part 3.} Fix \(p\in(0,1)\) and \(\eta>0\). Since \(\Phi\) is strictly
increasing, \(\delta=\min\{p-\Phi(z_p-\eta),\,\Phi(z_p+\eta)-p\}>0\). By
part~2 there is \(\nu_0\) with \(\sup_t|F_\nu-\Phi|<\delta\) for
\(\nu\geq\nu_0\); then \(F_\nu(z_p-\eta)<p<F_\nu(z_p+\eta)\). As \(F_\nu\)
is continuous and strictly increasing (\cref{thm:p15-student-law}) and
\(F_\nu(t_{\nu,p})=p\), this forces
\(z_p-\eta<t_{\nu,p}<z_p+\eta\) for all \(\nu\geq\nu_0\).
\end{proof}
\end{annexitem}

\begin{annexitem}{annex:p15-w02b-p15-059-mixture-order}{thm:p15-w02b-p15-059-mixture-order}{Scale mixture, strict ordering and undercoverage}
\begin{proof}
\emph{Mixture.} By \cref{thm:p15-student-law}, \(F_\nu\) is the CDF of
\(T=Z_0/\sqrt{W/\nu}\) with \(Z_0\sim N(0,1)\) and \(W\sim\chi^2_\nu\)
independent, and \(W>0\) almost surely by
\cref{thm:p15-w02b-p15-059-chisq-moments}. On \(\{W>0\}\),
\(\{T\leq t\}=\{Z_0\leq t\sqrt{W/\nu}\}\). The joint law is the product of
the two marginals, so Fubini's theorem gives
\(P(T\leq t)=\int P(Z_0\leq t\sqrt{w/\nu})\,dP_W(w)
=E\,\Phi(t\sqrt{W/\nu})=E\,G_t(W/\nu)\).

\emph{Strict concavity for \(t>0\).} On \((0,\infty)\),
\(G_t'(x)=\tfrac t2\,\phi_0(t\sqrt x)\,x^{-1/2}\) is a product of the
positive constant \(t/2\) and two positive strictly decreasing functions of
\(x\), so \(G_t'\) is strictly decreasing. For \(x>0\), \(x\neq1\), the mean
value theorem gives \(\xi\) strictly between \(1\) and \(x\) with
\(G_t(x)-G_t(1)=G_t'(\xi)(x-1)\). If \(x>1\) then \(G_t'(\xi)<G_t'(1)\) and
\(x-1>0\); if \(x<1\) then \(G_t'(\xi)>G_t'(1)\) and \(x-1<0\). In both
cases
\[
 G_t(x)<G_t(1)+G_t'(1)(x-1).
\]

\emph{Ordering.} Put \(X=W/\nu\). Then \(X>0\) almost surely, \(EX=1\)
by \cref{thm:p15-w02b-p15-059-chisq-moments}, and \(P(X=1)=0\) because the
law of \(W\) has no atoms. The random variable
\(D=G_t(1)+G_t'(1)(X-1)-G_t(X)\) is integrable (as \(G_t\) is bounded and
\(X\) integrable) and strictly positive almost surely, so \(ED>0\). Taking
expectations, \(G_t(1)-E\,G_t(X)=ED>0\), that is \(F_\nu(t)<\Phi(t)\) for
\(t>0\). For \(t<0\), symmetry of \(F_\nu\) and of \(\Phi\) gives
\(F_\nu(t)=1-F_\nu(-t)>1-\Phi(-t)=\Phi(t)\).

\emph{Quantiles.} If \(1/2<p<1\) then \(z_p>0\), so
\(F_\nu(z_p)<\Phi(z_p)=p=F_\nu(t_{\nu,p})\), and strict monotonicity of
\(F_\nu\) gives \(t_{\nu,p}>z_p\). If \(0<p<1/2\) then \(z_p<0\) and the
reverse inequality follows in the same way. Symmetry gives
\(F_\nu(0)=\Phi(0)=1/2\), so \(t_{\nu,1/2}=z_{1/2}=0\).

\emph{Coverage.} Under normal sampling \(V_n>0\) almost surely
(\cref{thm:p15-student-orthogonal}), and there
\(\mu\in C_c(X)\) if and only if \(|T_\mu|\leq c\). By
\cref{thm:p15-student-law} \(T_\mu\) has the atomless symmetric CDF
\(F_\nu\), so \(\kappa_\nu(c)=F_\nu(c)-F_\nu(-c)=2F_\nu(c)-1\), which is
less than \(2\Phi(c)-1\) because \(c>0\). With \(c=z_{(1+\alpha)/2}\),
\(2\Phi(c)-1=\alpha\). Finally \((1+\alpha)/2>1/2\), so
\(t_{\nu,(1+\alpha)/2}>z_{(1+\alpha)/2}\), and the Student half-width
\(t_{\nu,(1+\alpha)/2}e_n\) exceeds \(z_{(1+\alpha)/2}e_n\) whenever
\(e_n>0\), that is whenever \(V_n>0\).
\end{proof}
\end{annexitem}

\begin{annexitem}{annex:p15-w02b-p15-059-second-order}{thm:p15-w02b-p15-059-second-order}{The first correction term}
\begin{proof}
Fix a compact \(K\). Let \(W\sim\chi^2_\nu\), \(X=W/\nu\) and
\(\Delta=X-1\). By \cref{thm:p15-w02b-p15-059-chisq-moments},
\[
 E\Delta=0,\qquad \nu E\Delta^2=2,\qquad
 E\Delta^4=\frac{12\nu^2+48\nu}{\nu^4}\leq\frac{60}{\nu^2}\quad(\nu\geq1).
\]
Let \(A=\{X\geq1/2\}\). Chebyshev's inequality with the fourth moment gives
\(P(A^{\rm c})\leq P(|\Delta|>1/2)\leq16E\Delta^4\leq960/\nu^2\).

\emph{Constants.} Put \(M_K=\sup_{t\in K}|G_t'(1)|\leq
\phi_0(0)\sup_K|t|/2\). For \(y\geq1/2\),
\(|G_t''(y)|\leq\tfrac14|t|\phi_0(0)\,(y^{-3/2}+t^2y^{-1/2})
\leq\tfrac14|t|\phi_0(0)\,(2^{3/2}+\sqrt2\,t^2)\), so
\(B_K=\sup\{|G_t''(y)|:t\in K,\ y\geq1/2\}<\infty\). The function
\((t,y)\mapsto G_t''(y)\) is uniformly continuous on the compact set
\(K\times[1/2,3/2]\), so
\[
 \omega_K(\delta)=\sup\{|G_t''(y)-G_t''(1)|:t\in K,\ |y-1|\leq\delta\}
 \longrightarrow0\qquad(\delta\downarrow0).
\]

\emph{Decomposition.} By \cref{thm:p15-w02b-p15-059-mixture-order},
\(F_\nu(t)-\Phi(t)=E[G_t(X)-G_t(1)]=E[G_t(X)-G_t(1)-G_t'(1)\Delta]\),
using \(E\Delta=0\). Taylor's formula with integral remainder, valid
because \(G_t\) is twice continuously differentiable on \((0,\infty)\) and
the segment between \(1\) and \(X>0\) lies in \((0,\infty)\), gives
\[
 G_t(X)-G_t(1)-G_t'(1)\Delta
 =\Delta^2\int_0^1(1-s)\,G_t''(1+s\Delta)\,ds .
\]
Split \(E[\cdot]=E[\cdot\,\mathbf1_{A^{\rm c}}]+E[\cdot\,\mathbf1_A]\).

\emph{Off \(A\).} There \(0<X<1/2\), so \(|\Delta|<1\), and
\(|G_t(X)-G_t(1)-G_t'(1)\Delta|\leq1+M_K\). Hence
\(\nu\,|E[\cdots\mathbf1_{A^{\rm c}}]|\leq(1+M_K)\,960/\nu\to0\),
uniformly in \(t\in K\).

\emph{On \(A\).} Write the remainder as
\(\Delta^2\bigl[\tfrac12G_t''(1)+\rho_t\bigr]\) with
\(\rho_t=\int_0^1(1-s)\bigl(G_t''(1+s\Delta)-G_t''(1)\bigr)ds\).
First,
\[
 \nu E\bigl[\tfrac12G_t''(1)\Delta^2\mathbf1_A\bigr]
 =\tfrac12G_t''(1)\bigl(2-\nu E[\Delta^2\mathbf1_{A^{\rm c}}]\bigr),
\]
and \(\nu E[\Delta^2\mathbf1_{A^{\rm c}}]\leq\nu P(A^{\rm c})\leq960/\nu\)
because \(|\Delta|<1\) off \(A\). As \(|G_t''(1)|\leq B_K\), this term
equals \(G_t''(1)+O(1/\nu)\) uniformly on \(K\). Second, on \(A\) every
point \(1+s\Delta\), \(0\leq s\leq1\), lies between \(1\) and \(X\geq1/2\).
Fix \(\delta\in(0,1/2]\). On \(\{|\Delta|\leq\delta\}\),
\(|\rho_t|\leq\tfrac12\omega_K(\delta)\); on \(\{|\Delta|>\delta\}\cap A\),
\(|\rho_t|\leq B_K\). Hence, uniformly in \(t\in K\),
\[
\begin{aligned}
 \nu E\bigl[\Delta^2|\rho_t|\mathbf1_A\bigr]
 &\leq\tfrac12\omega_K(\delta)\,\nu E\Delta^2
   +B_K\,\nu E\bigl[\Delta^2\mathbf1_{\{|\Delta|>\delta\}}\bigr]\\
 &\leq\omega_K(\delta)+\frac{B_K\,\nu E\Delta^4}{\delta^2}
 \leq\omega_K(\delta)+\frac{60B_K}{\nu\delta^2},
\end{aligned}
\]
using \(\Delta^2\mathbf1_{\{|\Delta|>\delta\}}\leq\Delta^4/\delta^2\).

\emph{Conclusion.} Collecting the three pieces,
\[
 \limsup_{\nu\to\infty}\sup_{t\in K}
 \bigl|\nu(F_\nu(t)-\Phi(t))-G_t''(1)\bigr|\leq\omega_K(\delta)
\]
for every \(\delta\in(0,1/2]\), hence the left side is zero.
Since \(G_t''(1)=-\psi(t)\), the displayed uniform limit follows.

\emph{Part 1.} By \cref{thm:p15-w02b-p15-059-mixture-order},
\(\kappa_\nu(c)=2F_\nu(c)-1\), so
\(\nu[(2\Phi(c)-1)-\kappa_\nu(c)]=-2\nu(F_\nu(c)-\Phi(c))\to2\psi(c)
=c\phi_0(c)(1+c^2)/2\), taking \(K=\{c\}\).

\emph{Part 2.} Put \(t_\nu=t_{\nu,p}\) and \(K=[z_p-1,z_p+1]\). By
\cref{thm:p15-w02b-p15-059-student-limit}, \(t_\nu\in K\) for all large
\(\nu\) and \(t_\nu\to z_p\). Since \(F_\nu(t_\nu)=p=\Phi(z_p)\),
\[
 \nu\bigl(\Phi(t_\nu)-\Phi(z_p)\bigr)=-\nu\bigl(F_\nu(t_\nu)-\Phi(t_\nu)\bigr)
 =\psi(t_\nu)+r_\nu,
\]
where \(|r_\nu|\leq\sup_{t\in K}|\nu(F_\nu(t)-\Phi(t))+\psi(t)|\to0\).
Continuity of \(\psi\) gives
\(\nu(\Phi(t_\nu)-\Phi(z_p))\to\psi(z_p)\). By the mean value theorem
\(\Phi(t_\nu)-\Phi(z_p)=\phi_0(\eta_\nu)(t_\nu-z_p)\) with \(\eta_\nu\)
between \(t_\nu\) and \(z_p\), so \(\eta_\nu\to z_p\) and
\(\phi_0(\eta_\nu)\to\phi_0(z_p)>0\). Therefore
\(\nu(t_\nu-z_p)\to\psi(z_p)/\phi_0(z_p)=z_p(1+z_p^2)/4\).
\end{proof}
\end{annexitem}

\begin{annexitem}{annex:p15-w02b-p15-059-nonnormal}{thm:p15-w02b-p15-059-nonnormal}{Pointwise limits without normality}
\begin{proof}
\emph{Part 1.} For \(n\geq2\),
\[
 S_n^2=\frac n{n-1}\Bigl(\frac1n\sum_{i=1}^nX_i^2-\bar X_n^2\Bigr).
\]
The variables \(X_i\) and \(X_i^2\) are integrable, with means \(\mu\) and
\(\sigma^2+\mu^2\). By \cref{thm:p15-extended-mean-strong-law} applied to
each sequence, \(\bar X_n\to\mu\) and \(n^{-1}\sum_{i\leq n}X_i^2\to
\sigma^2+\mu^2\) almost surely, so \(S_n^2\to\sigma^2\) and, by continuity
of the square root, \(S_n\to\sigma\) almost surely. Almost sure convergence
implies convergence in probability.

\emph{Part 2.} By \cref{thm:p15-clt-iid},
\(\sqrt n(\bar X_n-\mu)=(\sum_{i\leq n}X_i-n\mu)/\sqrt n\Rightarrow
\sigma\chi\). Since \(S_n\to\sigma\neq0\) in probability, the ratio
statement of \cref{thm:p15-clt-perturbation} gives
\(\sqrt n(\bar X_n-\mu)/S_n\Rightarrow\sigma\chi/\sigma=\chi\), with any
fixed finite value on \(\{S_n=0\}\).

\emph{Part 3.} Since \(e_n=S_n/\sqrt n\), for every \(n\) and every
outcome, including those with \(V_n=0\),
\[
 \{\mu\in C_{c_n}(X)\}=\{Y_n\leq0\},\qquad
 Y_n=\bigl|\sqrt n(\bar X_n-\mu)\bigr|-c_nS_n .
\]
By \cref{thm:p15-weak-continuous-mapping} with the continuous map
\(x\mapsto|x|\), \(|\sqrt n(\bar X_n-\mu)|\Rightarrow\sigma|\chi|\). In
both choices \(c_n\to z=z_{(1+\alpha)/2}\) (by
\cref{thm:p15-w02b-p15-059-student-limit} for the Student choice), so
\(-c_nS_n\to-z\sigma\) in probability. The perturbation statement of
\cref{thm:p15-clt-perturbation} gives
\(Y_n\Rightarrow\sigma|\chi|-z\sigma\). This limit has a continuous CDF
because \(\sigma>0\) and \(\chi\) is atomless, so \(0\) is a continuity
point and
\[
 P(Y_n\leq0)\to P(|\chi|\leq z)=2\Phi(z)-1=\alpha .
\]
For the lower one-sided interval,
\(\{\bar X_n-c_n'e_n\leq\mu\}=\{\sqrt n(\bar X_n-\mu)-c_n'S_n\leq0\}\), and
the same two lemmas give
\(\sqrt n(\bar X_n-\mu)-c_n'S_n\Rightarrow\sigma\chi-z_\alpha\sigma\), so
the probability tends to \(\Phi(z_\alpha)=\alpha\); here
\(c_n'\to z_\alpha\). For the upper one,
\(\{\mu\leq\bar X_n+c_n'e_n\}=\{-\sqrt n(\bar X_n-\mu)-c_n'S_n\leq0\}\),
and \(-\sqrt n(\bar X_n-\mu)\Rightarrow-\sigma\chi\), which has the law of
\(\sigma\chi\); the limit is again \(\alpha\).

\emph{Part 4.} At \(m=\mu\), \(\phi_\mu^0=1\) exactly on
\(\{\mu\notin C_0(X)\}\), whose probability is
\(1-P\{\mu\in C_0(X)\}\to1-\alpha\) by part~3. For \(m\neq\mu\) the test
accepts exactly on \(\{|\sqrt n(\bar X_n-m)|\leq t_{\nu,(1+\alpha)/2}S_n\}\),
and \(|\sqrt n(\bar X_n-m)|\geq\sqrt n|\mu-m|-|\sqrt n(\bar X_n-\mu)|\), so
acceptance implies \(R_n\geq\sqrt n|\mu-m|\) with
\(R_n=|\sqrt n(\bar X_n-\mu)|+t_{\nu,(1+\alpha)/2}S_n\). As in part~3,
\(R_n\Rightarrow\sigma|\chi|+z\sigma\), which has a continuous CDF. Given
\(\eta>0\) choose \(M\) with \(P(\sigma|\chi|+z\sigma\geq M)<\eta\). For
\(n\) with \(\sqrt n|\mu-m|\geq M\), the acceptance probability is at most
\(P(R_n\geq M)\), which tends to \(P(\sigma|\chi|+z\sigma\geq M)<\eta\).
Hence the acceptance probability tends to zero.

\emph{Part 5.} Here \(\bar X_n-m=h\sigma/\sqrt n+\sigma\bar\xi_n\) and
\(V_n=\sigma^2\sum_{i\leq n}(\xi_i-\bar\xi_n)^2\), so \(S_n=\sigma S_n^\xi\),
where \(\bar\xi_n,S_n^\xi\) are computed from \(\xi_1,\ldots,\xi_n\). Parts~1
and~2 applied to the sequence \(\xi_i\) (mean \(0\), variance \(1\)) give
\(S_n\to\sigma\) almost surely and \(\sqrt n\bar\xi_n\Rightarrow\chi\);
hence \(\sqrt n(\bar X_n-m)=h\sigma+\sigma\sqrt n\bar\xi_n\Rightarrow
\sigma(h+\chi)\) by \cref{thm:p15-clt-perturbation}. The test rejects
exactly on \(\{Y_n>0\}\),
\(Y_n=|\sqrt n(\bar X_n-m)|-t_{\nu,(1+\alpha)/2}S_n\), and as in part~3
\(Y_n\Rightarrow\sigma|h+\chi|-z\sigma\), whose CDF is continuous. So the
rejection probability tends to
\(P(|h+\chi|>z)=P(\chi>z-h)+P(\chi<-z-h)=1-\Phi(z-h)+\Phi(-z-h)\).

\emph{Normal cutoffs.} Every step above used only \(c_n\to z\) (or
\(c_n'\to z_\alpha\)); the constant normal multipliers satisfy this, so
parts~3--5 hold for the normal-cutoff tests and intervals.
\end{proof}
\end{annexitem}

\begin{annexitem}{annex:p15-w02b-p15-059-nonuniform}{thm:p15-w02b-p15-059-nonuniform}{Bernoulli laws defeat every multiplier}
\begin{proof}
For \(0<p<1\) let \(P_p\) be the law with \(P_p\{1\}=p\),
\(P_p\{0\}=1-p\); its mean is \(p\) and its variance \(p(1-p)\in(0,1/4]\)
is finite and positive. On the event
\(E=\{X_1=\cdots=X_n=0\}\), of probability \((1-p)^n\), we have
\(\bar X_n=0\), \(V_n=0\) and \(e_n=0\), so
\(C_k(X)=\{0\}\) and \((-\infty,\bar X_n+ke_n]=(-\infty,0]\); neither
contains \(p>0\). By Bernoulli's inequality \((1-p)^n\geq1-np\), so both
coverages are at most \(1-(1-p)^n\leq np\), which tends to \(0\) as
\(p\downarrow0\).

For the lower one-sided interval use the reflected law \(P_p^-\) with
\(P_p^-\{-1\}=p\), \(P_p^-\{0\}=1-p\), of mean \(-p\) and the same
variance. On the same all-zero event the interval is \([0,\infty)\), which
does not contain \(-p\); its coverage is again at most \(np\).

Since coverage is nonnegative, each infimum is \(0\). For the final
assertion, the Student and normal-cutoff intervals are \(C_k\) for
particular \(k\geq0\); their infimum over the class is zero at every
\(n\geq2\), so it does not tend to \(\alpha\), while
\cref{thm:p15-w02b-p15-059-nonnormal} gives the pointwise limit \(\alpha\).
\end{proof}
\end{annexitem}

\begin{annexitem}{annex:p15-w02b-p15-059-uniform-regimes}{thm:p15-w02b-p15-059-uniform-regimes}{The uniform endpoint in two regimes}
\begin{proof}
\emph{Part 1.} By independence, \(P_\theta(L_n\leq x)=(x/\theta)^n\) for
\(0\leq x\leq\theta\); this CDF is continuous and \(L_n<\theta\) almost
surely. Since \((1-\alpha)^{-1/n}>1\), \(\theta\in I_n\) if and only if
\(L_n\leq\theta\leq L_n(1-\alpha)^{-1/n}\), that is (as \(L_n\leq\theta\)
almost surely) if and only if \(L_n\geq\theta(1-\alpha)^{1/n}\). Its
probability is \(1-\bigl((1-\alpha)^{1/n}\bigr)^n=\alpha\). For \(n=2\),
\((1-\alpha)^{-1/2}=1/q\) with \(q=\sqrt{1-\alpha}\), so
\(I_2=[L,L/q]=I_3^{\rm U}\).

\emph{Part 2.} \(|I_n|=L_n\bigl((1-\alpha)^{-1/n}-1\bigr)\) and
\(n\bigl((1-\alpha)^{-1/n}-1\bigr)=n(e^{a/n}-1)\to a\), since
\((e^{h}-1)/h\to1\) as \(h\to0\). The sequence \(L_n\) is nondecreasing and
bounded by \(\theta\), so it has an almost sure limit
\(L_\infty\leq\theta\); for \(0<\delta<\theta\),
\(P(L_\infty\leq\theta-\delta)\leq P(L_n\leq\theta-\delta)
=(1-\delta/\theta)^n\to0\). Hence \(L_n\to\theta\) almost surely and
\(n|I_n|\to a\theta\).

\emph{Part 3.} The uniform law on \((0,\theta)\) has mean \(\theta/2\) and
variance \(\theta^2/3-\theta^2/4=\theta^2/12>0\). By
\cref{thm:p15-clt-iid},
\(\sqrt n(\widehat\theta_n-\theta)=2\sqrt n(\bar X_n-\theta/2)\Rightarrow
2(\theta/\sqrt{12})\chi=(\theta/\sqrt3)\chi\), and by
\cref{thm:p15-extended-mean-strong-law}, \(\widehat\theta_n\to\theta\)
almost surely. Now \(\theta\in J_n\) if and only if \(Y_n\leq0\) with
\(Y_n=|\sqrt n(\widehat\theta_n-\theta)|-z\widehat\theta_n/\sqrt3\). By
\cref{thm:p15-weak-continuous-mapping,thm:p15-clt-perturbation},
\(Y_n\Rightarrow(\theta/\sqrt3)(|\chi|-z)\), whose CDF is continuous, so
\(P_\theta(Y_n\leq0)\to P(|\chi|\leq z)=\alpha\). Also
\(\sqrt n|J_n|=2z\widehat\theta_n/\sqrt3\to2z\theta/\sqrt3\) almost surely.

\emph{Part 4.} \(|I_n|/|J_n|=n|I_n|/\bigl(\sqrt n\cdot\sqrt n|J_n|\bigr)\),
whose numerator tends to \(a\theta\) and whose denominator tends to
infinity almost surely; so the ratio tends to zero. For the second
assertion write
\[
 \sqrt n\,(L_n-\min J_n)=\sqrt n\,(L_n-\theta)
   -\sqrt n\,(\widehat\theta_n-\theta)+\frac{z\widehat\theta_n}{\sqrt3}.
\]
For \(\delta>0\),
\(P\bigl(\sqrt n(\theta-L_n)>\delta\bigr)=(1-\delta/(\theta\sqrt n))^n
\leq\exp(-\delta\sqrt n/\theta)\to0\) once \(\delta<\theta\sqrt n\), so
the first term tends to zero in probability. The second tends weakly to
\(-(\theta/\sqrt3)\chi\), which has the law of \((\theta/\sqrt3)\chi\), and
the third tends to \(z\theta/\sqrt3\) almost surely. By
\cref{thm:p15-clt-perturbation} the sum tends weakly to
\((\theta/\sqrt3)(\chi+z)\), with continuous CDF. Hence
\(P_\theta(\min J_n<L_n)\to P(\chi>-z)=\Phi(z)=(1+\alpha)/2\).
\end{proof}
\end{annexitem}

\begin{annexitem}{annex:p15-w02b-p15-059-pvalue}{thm:p15-w02b-p15-059-pvalue}{The Student p-value and the significance event}
\begin{proof}
Put \(t=t_{\nu,(1+\alpha)/2}\); since \((1+\alpha)/2>1/2\),
\cref{thm:p15-w02b-p15-059-mixture-order} gives \(t>z_{(1+\alpha)/2}>0\).
By the definition of \(\phi_m^0\) in \cref{thm:p15-student-coverage},
\(\phi_m^0(X)=1\) exactly when \(|\bar X_n-m|>te_n\). As
\(C_0(X)=[\bar X_n-te_n,\bar X_n+te_n]\), also \(m\notin C_0(X)\) exactly
when \(|\bar X_n-m|>te_n\). It remains to show that \(\pi_m(X)<\varepsilon\)
exactly when \(|\bar X_n-m|>te_n\).

If \(V_n>0\), then \(e_n>0\) and \(|\bar X_n-m|>te_n\) means \(|T_m|>t\).
By \cref{thm:p15-student-law} \(F_\nu\) is continuous and strictly
increasing with \(F_\nu(t)=(1+\alpha)/2\), so \(|T_m|>t\) holds if and only
if \(F_\nu(|T_m|)>(1+\alpha)/2\), that is, if and only if
\(\pi_m(X)=2(1-F_\nu(|T_m|))<1-\alpha=\varepsilon\). If \(V_n=0\), then
\(e_n=0\) and \(|\bar X_n-m|>te_n\) means \(\bar X_n\neq m\), which is the
case \(\pi_m(X)=0<\varepsilon\); in the other case \(\pi_m(X)=1\), which is
not below \(\varepsilon<1\).

Now let \(\mu=m\) under normal sampling. Then \(V_n>0\) almost surely
(\cref{thm:p15-student-orthogonal}) and \(T_m\) has the CDF \(F_\nu\)
(\cref{thm:p15-student-law}), which is symmetric, continuous, and below
\(1\) everywhere (a strictly increasing CDF cannot reach \(1\)). For
\(s=0\), \(\{\pi_m\leq0\}=\{F_\nu(|T_m|)=1\}\) is empty on \(\{V_n>0\}\),
so its probability is \(0\). For \(s=1\) the probability is \(1\). For
\(0<s<1\) put \(q=t_{\nu,1-s/2}\), which is positive because
\(1-s/2>1/2\). Then \(\pi_m\leq s\) if and only if
\(F_\nu(|T_m|)\geq1-s/2=F_\nu(q)\), that is \(|T_m|\geq q\), and
\[
 P\{|T_m|\geq q\}=\bigl(1-F_\nu(q)\bigr)+F_\nu(-q)
 =2\bigl(1-F_\nu(q)\bigr)=s ,
\]
using symmetry and the absence of atoms. Hence
\(P_{m,\sigma}\{\pi_m\leq s\}=s\) for all \(s\in[0,1]\), which is the
inequality of \cref{def:p15-valid-p-value} with equality.
\end{proof}
\end{annexitem}

\begin{annexitem}{annex:p15-w02b-p15-059-verdicts}{thm:p15-w02b-p15-059-verdicts}{Verdicts at a fixed sample size}
\begin{proof}
Write \(L=\bar X_n-te_n\leq U=\bar X_n+te_n\), so \(C_0(X)=[L,U]\), with
\(t=t_{\nu,(1+\alpha)/2}\).

\emph{Two descriptions of the verdicts.} First, \([L,U]\subset
(m-\delta,m+\delta)\) if and only if \(m-\delta<L\) and \(U<m+\delta\).
Second, \([L,U]\cap(m-\delta,m+\delta)=\emptyset\) if and only if
\(L\geq m+\delta\) or \(U\leq m-\delta\). In the second statement the
backward implication is immediate. For the forward one suppose
\(L<m+\delta\) and \(U>m-\delta\). If \(L>m-\delta\), the point \(L\)
lies in both sets. If \(L\leq m-\delta\), put
\(b=\min\{U,m+\delta\}>m-\delta\) and \(x=(m-\delta+b)/2\); then
\(L\leq m-\delta<x\leq U\) and \(m-\delta<x<m+\delta\), so \(x\) lies in
both sets.

\emph{Part 1.} On \(\mathcal O_n\) the interval misses
\(B_\delta\ni m\), so \(m\notin C_0(X)\) and \(\phi_m^0(X)=1\) by
\cref{thm:p15-w02b-p15-059-pvalue}. On \(\mathcal N_n\cap\{\phi_m^0=1\}\),
\(m\notin[L,U]\), so either \(m<L\), and then \([L,U]\subset(m,m+\delta)\)
because \(U<m+\delta\), or \(U<m\), and then \([L,U]\subset(m-\delta,m)\)
because \(L>m-\delta\). Conversely, if \([L,U]\) lies in
\((m,m+\delta)\) or in \((m-\delta,m)\), it lies in \(B_\delta\) and
misses \(m\), so the outcome is in \(\mathcal N_n\cap\{\phi_m^0=1\}\).

\emph{Part 2.} With coefficient \((1+\alpha)/2\) the one-sided critical
value of \cref{thm:p15-student-coverage} is \(t\). The test
\(\phi_{m-\delta}^+\) rejects when \(\bar X_n-(m-\delta)>te_n\), that is
when \(L>m-\delta\); the test \(\phi_{m+\delta}^-\) rejects when
\(\bar X_n-(m+\delta)<-te_n\), that is when \(U<m+\delta\). By the first
description, both reject exactly on \(\mathcal N_n\).

\emph{Part 3.} Under normal sampling \(V_n>0\) almost surely and
\(T_\mu\) has the symmetric atomless CDF \(F_\nu\)
(\cref{thm:p15-student-orthogonal,thm:p15-student-law}). If
\(\mu-m\geq\delta\), then on \(\mathcal N_n\cap\{V_n>0\}\),
\(U<m+\delta\leq\mu\), so \(\bar X_n-\mu<-te_n\), that is \(T_\mu<-t\);
hence \(P_{\mu,\sigma}(\mathcal N_n)\leq F_\nu(-t)=1-F_\nu(t)=(1-\alpha)/2\).
If \(\mu-m\leq-\delta\), then on \(\mathcal N_n\cap\{V_n>0\}\),
\(L>m-\delta\geq\mu\), so \(T_\mu>t\), with the same bound.

\emph{Part 4.} Let \(m-\delta\leq\mu\leq m+\delta\). By the second
description, on \(\mathcal O_n\cap\{V_n>0\}\) either
\(L\geq m+\delta\geq\mu\), so \(T_\mu\geq t\), or
\(U\leq m-\delta\leq\mu\), so \(T_\mu\leq-t\). Hence
\(P_{\mu,\sigma}(\mathcal O_n)\leq P\{|T_\mu|\geq t\}=2(1-F_\nu(t))=1-\alpha\).

\emph{Part 5.} \(C_0\) has coverage exactly \(\alpha\) at every
\((\mu,\sigma)\) by \cref{thm:p15-student-coverage}; at \(\mu=m\),
\(P_{m,\sigma}\{\phi_m^0=1\}=P_{m,\sigma}\{m\notin C_0(X)\}=1-\alpha\) by
\cref{thm:p15-w02b-p15-059-pvalue}.
\end{proof}
\end{annexitem}

\begin{annexitem}{annex:p15-w02b-p15-059-scale}{thm:p15-w02b-p15-059-scale}{Significance responds to the departure in units of the scale}
\begin{proof}
\emph{Part 1.} Under \(P_{\mu,\sigma}\) put \(Z_i=(X_i-\mu)/\sigma\); the
\(Z_i\) are independent standard normal, a law not depending on
\((\mu,\sigma)\). Let \(\bar Z\), \(V^Z\) and \(e^Z\) be computed from
\(Z_1,\ldots,Z_n\) as in \cref{def:p15-student-model}. Then, with
\(\vartheta=(\mu-m)/\sigma\), \(\bar X_n-m=\sigma(\vartheta+\bar Z)\) and
\(V_n=\sigma^2V^Z\), so \(e_n=\sigma e^Z\). Since \(\sigma>0\),
\(\phi_m^0(X)=1\) exactly when \(|\vartheta+\bar Z|>te^Z\). Put
\(\beta_n(\vartheta)=P\{|\vartheta+\bar Z|>te^Z\}\), computed under the
standard normal sample; this is the rejection probability at every
\((\mu,\sigma)\) with \((\mu-m)/\sigma=\vartheta\).

\emph{Part 2.} \(\beta_n(0)=P_{0,1}\{\phi_0^0=1\}=1-\alpha\) by
\cref{thm:p15-w02b-p15-059-verdicts}, part~5. For continuity, first
\(P\{|\vartheta+\bar Z|=te^Z\}=0\) for every \(\vartheta\): by
\cref{thm:p15-student-orthogonal} with \(\mu=0\), \(\sigma=1\), the
variable \(Y_1=\sqrt n\bar Z\) is standard normal and independent of
\(W=V^Z\), and \(e^Z=\sqrt{W/(n\nu)}\). For each fixed \(w\geq0\) the set
of \(y\) with \(|\vartheta+y/\sqrt n|=t\sqrt{w/(n\nu)}\) has at most two
points, so has probability zero under the law of \(Y_1\); integrating over
the law of \(W\) (Fubini's theorem for the product law) gives zero. Now
let \(\vartheta_k\to\vartheta\). Off this null event the indicators
\(\mathbf1\{|\vartheta_k+\bar Z|>te^Z\}\) converge to
\(\mathbf1\{|\vartheta+\bar Z|>te^Z\}\), because the inequality is strict
in one direction or the other at \(\vartheta\) and both sides are
continuous in \(\vartheta\). Bounded convergence gives
\(\beta_n(\vartheta_k)\to\beta_n(\vartheta)\).

For the limit, let \(\vartheta_k\to+\infty\). Then
\(1-\beta_n(\vartheta_k)=P\{|\vartheta_k+\bar Z|\leq te^Z\}\leq
P\{\vartheta_k+\bar Z\leq te^Z\}\), and the indicator of
\(\{\vartheta_k+\bar Z\leq te^Z\}\) tends to zero at every outcome because
\(\bar Z\) and \(e^Z\) are finite; bounded convergence gives
\(\beta_n(\vartheta_k)\to1\). For \(\vartheta_k\to-\infty\) use
\(\{|\vartheta_k+\bar Z|\leq te^Z\}\subset
\{-\vartheta_k-\bar Z\leq te^Z\}\) in the same way. As the sequence was
arbitrary, \(\beta_n(\vartheta)\to1\) as \(|\vartheta|\to\infty\).

\emph{Part 3.} By continuity at \(0\) there is \(r>0\) with
\(\beta_n(\vartheta)<1-\alpha+\eta\) for \(|\vartheta|<r\); choose
\(\vartheta_1\in(0,r)\) and \(\sigma_1=D/\vartheta_1\). By the limit there
is \(\vartheta_2>0\) with \(\beta_n(\vartheta_2)>1-\eta\); put
\(\sigma_2=D/\vartheta_2\). With \(\mu=m+D\),
\((\mu-m)/\sigma_j=\vartheta_j\), and part~1 gives the two inequalities.
For the final sentence take \(0<D<D'\) and \(0<\eta<\alpha/2\). Applying
the construction to \(D'\) gives a scale at which the larger departure
\(D'\) has rejection probability below \(1-\alpha+\eta\), and applying it
to \(D\) gives a scale at which the smaller departure \(D\) has rejection
probability above \(1-\eta>1-\alpha+\eta\).
\end{proof}
\end{annexitem}

\begin{annexitem}{annex:p15-w02b-p15-059-magnitude-limits}{thm:p15-w02b-p15-059-magnitude-limits}{Significance against magnitude in large samples}
\begin{proof}
Write \(t_\nu=t_{\nu,(1+\alpha)/2}\), \(L_n=\bar X_n-t_\nu e_n\) and
\(U_n=\bar X_n+t_\nu e_n\). By
\cref{thm:p15-w02b-p15-059-nonnormal}, part~1, \(S_n\to\sigma\) almost
surely; by \cref{thm:p15-extended-mean-strong-law}, \(\bar X_n\to\mu\)
almost surely; and \(t_\nu\to z\) by
\cref{thm:p15-w02b-p15-059-student-limit}. Let \(\Omega_0\) be the event of
probability one on which both almost sure limits hold.

\emph{Part 1.} The two statements about \(P\{\phi_m^0=1\}\) are part~4 of
\cref{thm:p15-w02b-p15-059-nonnormal}. On \(\Omega_0\),
\(S_n>0\) for all large \(n\), so \(V_n=(n-1)S_n^2>0\), and then
\(T_m/\sqrt n=(\bar X_n-m)/S_n\to d/\sigma\). If \(d\neq0\) this gives
\(|T_m|\to\infty\) on \(\Omega_0\), and for large \(n\)
\[
 0\leq\pi_m(X)=2\bigl(1-F_\nu(|T_m|)\bigr)
 \leq2\bigl(1-\Phi(|T_m|)\bigr)+2\sup_{x\in\mathbb R}|F_\nu(x)-\Phi(x)| .
\]
The first term tends to zero because \(|T_m|\to\infty\), and the second by
\cref{thm:p15-w02b-p15-059-student-limit}, part~2.

\emph{Part 2.} On \(\Omega_0\), \(t_\nu e_n=t_\nu S_n/\sqrt n\to0\), so
\(L_n\to\mu\) and \(U_n\to\mu\); and
\(\sqrt n\,|C_0(X)|=2t_\nu S_n\to2z\sigma\).

\emph{Part 3.} The two descriptions in the proof of
\cref{thm:p15-w02b-p15-059-verdicts} hold for every outcome:
\(\mathcal N_n=\{m-\delta<L_n\}\cap\{U_n<m+\delta\}\) and
\(\mathcal O_n=\{L_n\geq m+\delta\}\cup\{U_n\leq m-\delta\}\).

If \(|d|<\delta\), then \(\mu\in B_\delta\), an open set, so on
\(\Omega_0\) both \(L_n\) and \(U_n\) lie in \(B_\delta\) for all large
\(n\), and the outcome is in \(\mathcal N_n\). Bounded convergence gives
\(P(\mathcal N_n)\to1\), and \(P(\mathcal O_n)\leq1-P(\mathcal N_n)\to0\)
because the verdicts are disjoint. If \(d>\delta\), then
\(\mu>m+\delta\), so on \(\Omega_0\) eventually \(L_n>m+\delta\) and the
outcome is in \(\mathcal O_n\); if \(d<-\delta\), eventually
\(U_n<m-\delta\). In both cases \(P(\mathcal O_n)\to1\) and
\(P(\mathcal N_n)\to0\).

Let \(d=\delta\), so \(\mu=m+\delta\). Put
\(Y_n^\pm=\sqrt n(\bar X_n-\mu)\pm t_\nu S_n\). By \cref{thm:p15-clt-iid},
\(\sqrt n(\bar X_n-\mu)\Rightarrow\sigma\chi\); since \(t_\nu S_n\to
z\sigma\) almost surely, hence in probability,
\cref{thm:p15-clt-perturbation} gives \(Y_n^\pm\Rightarrow\sigma(\chi\pm
z)\), whose CDF is continuous because \(\sigma>0\). For a sequence
\(Y_n\Rightarrow Y\) with continuous CDF \(G\),
\(P(Y_n\leq-\eta)\leq P(Y_n<0)\leq P(Y_n\leq0)\) for \(\eta>0\), and
letting \(n\to\infty\) and then \(\eta\downarrow0\) gives
\(P(Y_n<0)\to G(0)\); likewise \(P(Y_n\geq0)\to1-G(0)\). Now
\(\{U_n<m+\delta\}=\{Y_n^+<0\}\) has probability tending to
\(P\{\chi<-z\}=1-\Phi(z)=(1-\alpha)/2\), while
\(P\{m-\delta<L_n\}\to1\) because \(L_n\to m+\delta\) on \(\Omega_0\).
Since \(|P(A\cap B)-P(A)|\leq P(B^{\rm c})\), \(P(\mathcal N_n)\to
(1-\alpha)/2\). Next \(\{L_n\geq m+\delta\}=\{Y_n^-\geq0\}\) has
probability tending to \(P\{\chi\geq z\}=(1-\alpha)/2\), and
\(P\{U_n\leq m-\delta\}\to0\) because \(U_n\to m+\delta\) on
\(\Omega_0\). Since \(P(A)\leq P(A\cup B)\leq P(A)+P(B)\),
\(P(\mathcal O_n)\to(1-\alpha)/2\).

Let \(d=-\delta\). The variables \(X_i'=2m-X_i\) are independent with a
common law of mean \(2m-\mu=m+\delta\) and variance \(\sigma^2\). Their
sample has \(\bar X_n'=2m-\bar X_n\) and the same \(S_n\), so
\(C_0(X')=\{2m-x:x\in C_0(X)\}\); reflection in \(m\) maps \(B_\delta\)
onto itself, so the verdict events of \(X'\) and of \(X\) coincide, and the
case \(d=\delta\) applies.

\emph{Part 4.} If \(0<d<\delta\), then \(\mu\) lies in the open interval
\((m,m+\delta)\), so on \(\Omega_0\) eventually \(L_n\) and \(U_n\) lie in
it, hence \(C_0(X)\subset(m,m+\delta)\); then \(m\notin C_0(X)\) and the
test rejects (\cref{thm:p15-w02b-p15-059-pvalue}). Bounded convergence
gives probability tending to \(1\). The case \(-\delta<d<0\) is the same
with \((m-\delta,m)\).
\end{proof}
\end{annexitem}

\begin{annexitem}{annex:p15-w02b-p15-059-posterior}{thm:p15-w02b-p15-059-posterior}{Posterior statements in large samples}
\begin{proof}
\emph{Part 1.} Since \(\tau^{-2}+n\sigma^{-2}=(n+\sigma^2\tau^{-2})/\sigma^2\),
\(v_n=\sigma^2/(n+\sigma^2\tau^{-2})\) and
\(nv_n=\sigma^2n/(n+\sigma^2\tau^{-2})\to\sigma^2\). Also
\(nv_n\sigma^{-2}-1=-\sigma^2\tau^{-2}/(n+\sigma^2\tau^{-2})
=-v_n\tau^{-2}\), so
\[
 u_n-\bar X_n=v_n\mu_0\tau^{-2}+\bar X_n\bigl(nv_n\sigma^{-2}-1\bigr)
 =v_n\tau^{-2}(\mu_0-\bar X_n).
\]
By \cref{thm:p15-extended-mean-strong-law}, \(\bar X_n\to\mu\) almost
surely; as \(\sqrt n\,v_n=nv_n/\sqrt n\to0\),
\(\sqrt n(u_n-\bar X_n)\to0\) almost surely, and so \(u_n\to\mu\).

\emph{Part 2.} Fix an outcome in the almost sure event of part~1. If \(G\)
is open with \(\mu\in G\), choose \(r>0\) with \((\mu-2r,\mu+2r)\subset
G\); if \(F\) is closed with \(\mu\notin F\), its complement is open and
contains \(\mu\), so choose \(r>0\) with \((\mu-2r,\mu+2r)\cap
F=\emptyset\). In both cases it suffices that
\(\Pi_n(|\theta-\mu|\geq2r)\to0\). For \(n\) with \(|u_n-\mu|<r\),
\(\{|\theta-\mu|\geq2r\}\subset\{|\theta-u_n|\geq r\}\), and Chebyshev's
inequality for the law \(\operatorname{normal}(u_n,v_n)\) gives
\(\Pi_n(|\theta-u_n|\geq r)\leq v_n/r^2\to0\).

\emph{Part 3.} Apply part~2 with \(G=(m,\infty)\) when \(d>0\);
\(G=(m-\delta,m+\delta)\) when \(|d|<\delta\);
\(G=(m,m+\delta)\) when \(0<d<\delta\); and the closed set
\(F=[m-\delta,m+\delta]\), which does not contain \(\mu\) when
\(|d|>\delta\), together with \(\Pi_n(|\theta-m|<\delta)\leq\Pi_n(F)\).

\emph{Part 4.} By symmetry of \(\Phi\),
\(\Pi_n(\theta>m)=1-\Phi\bigl((m-u_n)/\sqrt{v_n}\bigr)
=\Phi\bigl((u_n-m)/\sqrt{v_n}\bigr)\). With \(\mu=m\),
\[
 (u_n-m)/\sqrt{v_n}=A_nY_n+B_n,
\]
where \(Y_n=\sqrt n(\bar X_n-\mu)/\sigma\), \(A_n=\sigma/\sqrt{nv_n}\) and
\(B_n=(u_n-\bar X_n)/\sqrt{v_n}=\sqrt{v_n}\,\tau^{-2}(\mu_0-\bar X_n)\).
By \cref{thm:p15-student-orthogonal}, \(Y_n\) is standard normal for every
\(n\), so \(Y_n\Rightarrow\chi\); \(A_n\to1\) by part~1; and \(B_n\to0\)
almost surely because \(v_n\to0\) and \(\bar X_n\to\mu\). Hence
\((u_n-m)/\sqrt{v_n}\Rightarrow\chi\) by
\cref{thm:p15-clt-perturbation}, and since \(\Phi\) is continuous,
\cref{thm:p15-weak-continuous-mapping} gives
\(\Pi_n(\theta>m)\Rightarrow\Phi(\chi)\). For \(0<s<1\),
\(P\{\Phi(\chi)\leq s\}=P\{\chi\leq z_s\}=s\), so \(\Phi(\chi)\) is
uniform on \((0,1)\).

\emph{Part 5.} \(\Pi_n(K_n)=\Phi(z)-\Phi(-z)=2\Phi(z)-1=\alpha\), and
\(\sqrt n\,|K_n|=2z\sqrt{nv_n}\to2z\sigma\) by part~1. The endpoint
differences are \((u_n-\bar X_n)\pm z(\sqrt{v_n}-\sigma/\sqrt n)\); after
multiplication by \(\sqrt n\) they become
\(\sqrt n(u_n-\bar X_n)\pm z(\sqrt{nv_n}-\sigma)\), which tend to zero
almost surely by part~1.
\end{proof}
\end{annexitem}
